% !TEX program = xelatex
% Lernplatt - by Can Siebert
% Kurs 22 (Bo 143) - Tag 9 - Aufgabenheft
% Performance-Marketing: Anzeigen, Conversion & Landing Pages

\documentclass[11pt,a4paper,oneside]{article}

\usepackage{polyglossia}
\setdefaultlanguage{german}

\usepackage[a4paper,margin=2.5cm]{geometry}

\usepackage{fontspec}
\defaultfontfeatures{Ligatures=TeX}

\IfFontExistsTF{Charter}{\setmainfont{Charter}}{%
  \IfFontExistsTF{Bitstream Charter}{\setmainfont{Bitstream Charter}}{%
    \setmainfont{TeX Gyre Schola}}}
\IfFontExistsTF{Source Sans 3}{\setsansfont{Source Sans 3}}{%
  \IfFontExistsTF{Source Sans Pro}{\setsansfont{Source Sans Pro}}{%
    \setsansfont{TeX Gyre Heros}}}

\newcommand{\caveatfontfallback}{\itshape}
\IfFontExistsTF{Caveat}{\newfontfamily\caveatfont{Caveat}[Scale=1.15]}{\newcommand\caveatfont{\caveatfontfallback}}

\setmonofont{Latin Modern Mono}

\usepackage{microtype}
\linespread{1.15}

\usepackage{xcolor}
\definecolor{lpInk}{RGB}{20,28,38}
\definecolor{lpAccent}{RGB}{22,90,140}
\definecolor{lpAccentLight}{RGB}{210,228,240}
\definecolor{lpAmber}{RGB}{222,138,30}
\definecolor{lpAmberLight}{RGB}{255,243,217}
\definecolor{lpAmberBorder}{RGB}{196,118,18}
\definecolor{lpGray}{RGB}{96,104,114}
\definecolor{lpGrayLight}{RGB}{238,240,243}
\definecolor{lpGreen}{RGB}{34,120,72}
\definecolor{lpGreenLight}{RGB}{222,238,228}
\definecolor{lpGreenBorder}{RGB}{24,96,58}
\definecolor{lpL1}{RGB}{34,120,72}
\definecolor{lpL2}{RGB}{22,90,140}
\definecolor{lpL3}{RGB}{170,42,52}

\usepackage[most]{tcolorbox}
\tcbuselibrary{breakable,skins}

\newtcolorbox{infobox}[1][Hinweis]{enhanced, breakable, colback=lpAccentLight, colframe=lpAccent, boxrule=0.6pt, arc=2pt, left=10pt,right=10pt,top=8pt,bottom=8pt, fonttitle=\sffamily\bfseries\color{white}, coltitle=white, title={#1}, attach boxed title to top left={xshift=8pt,yshift=-8pt}, boxed title style={size=small, colback=lpAccent, colframe=lpAccent, boxrule=0.4pt, arc=1pt}, top=14pt}

\newtcolorbox{antwortbox}[1][Ihre Antwort]{enhanced, breakable, colback=white, colframe=lpGray, boxrule=0.4pt, arc=2pt, left=10pt,right=10pt,top=8pt,bottom=8pt, fonttitle=\sffamily\bfseries\color{white}, coltitle=white, title={#1}, attach boxed title to top left={xshift=8pt,yshift=-8pt}, boxed title style={size=small, colback=lpGray, colframe=lpGray, boxrule=0.4pt, arc=1pt}, top=14pt}

\usepackage{enumitem}
\setlist{itemsep=2pt, topsep=4pt, parsep=0pt}
\setlist[itemize,1]{label={\textbullet},leftmargin=1.6em}
\setlist[itemize,2]{label={\textendash},leftmargin=1.4em}
\setlist[enumerate,1]{label=\arabic*., leftmargin=1.8em}

\usepackage{tabularx}
\usepackage{array}
\usepackage{booktabs}
\usepackage{longtable}

\usepackage{fancyhdr}
\usepackage{lastpage}
\pagestyle{fancy}
\fancyhf{}
\renewcommand{\headrulewidth}{0.3pt}
\renewcommand{\footrulewidth}{0pt}
\fancyhead[L]{\footnotesize\sffamily\color{lpGray}Aufgabenheft \textperiodcentered{} Kurs 22 \textperiodcentered{} Tag 9}
\fancyhead[R]{\footnotesize\sffamily\color{lpGray}Performance-Marketing}
\fancyfoot[L]{\footnotesize\sffamily\color{lpGray}\textcopyright{} Can Siebert}
\fancyfoot[R]{\footnotesize\sffamily\color{lpGray}\thepage{} / \pageref{LastPage}}

\fancypagestyle{plain}{\fancyhf{}\renewcommand{\headrulewidth}{0pt}\fancyfoot[C]{\footnotesize\sffamily\color{lpGray}\thepage{} / \pageref{LastPage}}}

\usepackage[explicit]{titlesec}
\titleformat{\section}[block]{\sffamily\Large\bfseries\color{lpAccent}}{\thesection.}{0.6em}{#1}[{\vspace{2pt}\color{lpAccent}\titlerule[0.6pt]}]
\titleformat{\subsection}[block]{\sffamily\large\bfseries\color{lpInk}}{\thesubsection}{0.6em}{#1}
\titlespacing*{\section}{0pt}{14pt}{8pt}
\titlespacing*{\subsection}{0pt}{10pt}{4pt}

\usepackage[hidelinks,unicode]{hyperref}

\newcommand{\akzent}[1]{{\caveatfont\color{lpAmber}#1}}
\newcommand{\fachbegriff}[1]{\textbf{#1}}

\newcommand{\niveaubadge}[2]{\tcbox[on line, colback=#1, colframe=#1, coltext=white, boxrule=0pt, arc=2pt, left=5pt,right=5pt,top=1pt,bottom=1pt, nobeforeafter]{\sffamily\bfseries\footnotesize #2}}
\newcommand{\nivLi}{\niveaubadge{lpL1}{L1\,\textbullet\,Wissen}}
\newcommand{\nivLii}{\niveaubadge{lpL2}{L2\,\textbullet\,Anwendung}}
\newcommand{\nivLiii}{\niveaubadge{lpL3}{L3\,\textbullet\,Management}}

\newcommand{\zeit}[1]{\tcbox[on line, colback=lpGrayLight, colframe=lpGrayLight, coltext=lpInk, boxrule=0pt, arc=2pt, left=5pt,right=5pt,top=1pt,bottom=1pt, nobeforeafter]{\sffamily\footnotesize\textbf{$\sim$\,#1}}}

\newcommand{\aufgabenkopf}[4]{\par\vspace{6pt}\noindent{\sffamily\Large\bfseries\color{lpAccent}Aufgabe #1 \textendash{} #2}\par\vspace{2pt}\noindent{\color{lpAccent}\rule{\linewidth}{0.6pt}}\par\vspace{4pt}\noindent #3 \hfill \zeit{#4}\par\vspace{6pt}}

\newcommand{\ytquery}[1]{\par\vspace{4pt}\noindent{\footnotesize\sffamily\color{lpGray}\textbf{YouTube-Suchquery:} \texttt{#1}}\par}

\newcommand{\antwortlinien}[1]{\par\vspace{6pt}\foreach \x in {1,...,#1}{\noindent\makebox[\linewidth]{\dotfill}\\[6pt]}}
\usepackage{pgffor}

\begin{document}

\begin{titlepage}
  \pagestyle{empty}
  \vspace*{1.5cm}
  \noindent{\sffamily\footnotesize\color{lpGray}LERNPLATT \textperiodcentered{} BY CAN SIEBERT}\\[2pt]
  \noindent{\color{lpAccent}\rule{\linewidth}{1.2pt}}\\[4pt]
  \noindent{\sffamily\footnotesize\color{lpGray}AUFGABENHEFT \textperiodcentered{} MODUL 5 \textperiodcentered{} TAG 9 \textperiodcentered{} KURS 22 (Bo 143) \textperiodcentered{} 10.\,Juni 2026}

  \vspace{3cm}
  {\sffamily\fontsize{30}{34}\selectfont\bfseries\color{lpInk}Aufgabenheft\\Tag 9\par}

  \vspace{0.7cm}
  {\sffamily\large\color{lpGray}Performance-Marketing:\\Anzeigen, Conversion \& Landing Pages\\\textendash{} elf Aufgaben auf drei Niveaus.\par}

  \vspace{1.5cm}
  {\caveatfont\LARGE\color{lpAmber}11 Aufgaben \textperiodcentered{} L1 \textperiodcentered{} L2 \textperiodcentered{} L3}\par

  \vfill

  \noindent\begin{minipage}{0.55\linewidth}
    {\sffamily\footnotesize\color{lpGray}Niveau-Mix:}\\
    {\sffamily\small\color{lpInk}3\,\textsf{L1} \textperiodcentered{} 5\,\textsf{L2} \textperiodcentered{} 3\,\textsf{L3}}\\[10pt]
    {\sffamily\footnotesize\color{lpGray}Bearbeitung:}\\
    {\sffamily\small\color{lpInk}Selbstbearbeitung im Lernpfad}
  \end{minipage}\hfill
  \begin{minipage}{0.4\linewidth}
    \raggedleft
    {\sffamily\footnotesize\color{lpGray}Dozent:}\\
    {\sffamily\small\color{lpInk}Can Siebert}\\[10pt]
    {\sffamily\footnotesize\color{lpGray}Begleitend:}\\
    {\sffamily\small\color{lpInk}Lehrskript \& Lösungsheft}
  \end{minipage}

  \vspace{0.5cm}
  \noindent{\color{lpAccent}\rule{\linewidth}{0.6pt}}
\end{titlepage}

\section*{So arbeiten Sie mit diesem Heft}
\thispagestyle{plain}

\noindent Tag 9 eröffnet Modul 5 \textendash{} \emph{Performance-Marketing als messbarer Vertriebshebel}. Anzeigen sind nicht isolierte Werbemaßnahme, sondern Teil eines Vertriebssystems aus Zielgruppenanalyse, Kampagnenplanung, Landing Pages, Conversion-Optimierung und Ergebnissteuerung. Drei Begriffe heute: \fachbegriff{Funnel}, \fachbegriff{Landing Page} und \fachbegriff{Leadqualität}.

\begin{itemize}
  \item \nivLi{} \textendash{} \fachbegriff{Wissen.} Funnel-Logik, KPIs, Anzeigenkanäle, Landing-Page-Elemente.
  \item \nivLii{} \textendash{} \fachbegriff{Anwendung.} SolarOffice GmbH und DataSecure Solutions als Cases.
  \item \nivLiii{} \textendash{} \fachbegriff{Management.} Performance-Marketing-Architektur als Investitionsentscheidung.
\end{itemize}

\begin{infobox}[Hinweis]
Performance-Marketing endet nicht bei Klicks. Es endet beim \emph{Vertriebsergebnis} \textendash{} Pipeline, Umsatz, Deckungsbeitrag. Wer Performance-Marketing nur auf Kosten pro Klick optimiert, kauft Aktivität, nicht Wirkung.
\end{infobox}

\clearpage

\section*{Niveau L1 \textendash{} Wissen \& Reproduktion}
\addcontentsline{toc}{section}{L1 -- Wissen}

% --- AUFGABE 1 -----------------------------------------------------------
% LERNPLATT_HILFSVIDEOS
\section*{Hilfsvideos zu diesem Tag}
\addcontentsline{toc}{section}{Hilfsvideos zu diesem Tag}
\thispagestyle{plain}

\noindent Die folgenden YouTube-Erkl\"arvideos vermitteln die Inhalte, die Sie zur Bearbeitung der Aufgaben ben\"otigen. Klicken Sie auf den Titel, um das Video direkt zu \"offnen; die vollst\"andige URL steht in Klammern.

\begin{description}[style=nextline,leftmargin=18pt,labelindent=0pt,itemsep=8pt]
  \item[1.~Performance-Marketing als messbarer Vertriebshebel] \href{https://youtu.be/756dea00Wmg}{\textcolor{lpAccent}{https://youtu.be/756dea00Wmg}}\\[2pt]
  {\footnotesize\color{lpGray}(URL: \nolinkurl{https://youtu.be/756dea00Wmg})}
  \item[2.~Landing Page Optimierung: Conversion-Hebel in der Praxis] \href{https://youtu.be/ZOE-CCudzZM}{\textcolor{lpAccent}{https://youtu.be/ZOE-CCudzZM}}\\[2pt]
  {\footnotesize\color{lpGray}(URL: \nolinkurl{https://youtu.be/ZOE-CCudzZM})}
  \item[3.~Sales Funnel: vom Klick zum Abschluss] \href{https://youtu.be/ZBEnOEcXtcs}{\textcolor{lpAccent}{https://youtu.be/ZBEnOEcXtcs}}\\[2pt]
  {\footnotesize\color{lpGray}(URL: \nolinkurl{https://youtu.be/ZBEnOEcXtcs})}
  \item[4.~Kennzahlen lesen: KPI-Architektur entlang des Funnels] \href{https://youtu.be/hmpDH5K5jRU}{\textcolor{lpAccent}{https://youtu.be/hmpDH5K5jRU}}\\[2pt]
  {\footnotesize\color{lpGray}(URL: \nolinkurl{https://youtu.be/hmpDH5K5jRU})}
\end{description}
\vspace{0.6em}
\noindent{\footnotesize\color{lpGray}\textit{Tipp:} \"Offnen Sie das Video, lassen Sie es im Hintergrund laufen und arbeiten Sie parallel an der Aufgabe.}

\clearpage

\aufgabenkopf{1}{Reichweite \textendash{} Traffic \textendash{} Leads \textendash{} Conversion \textendash{} Umsatz}{\nivLi}{6\,min}

\noindent\textbf{Aufgabe:}

\begin{enumerate}
  \item Definieren Sie in einem Satz: \emph{Reichweite \textperiodcentered{} Traffic \textperiodcentered{} Leads \textperiodcentered{} Conversion \textperiodcentered{} Umsatz}.
  \item Auf welcher Ebene wird Performance-Marketing im Mittelstand am häufigsten \emph{gemessen} \textendash{} und auf welcher Ebene müsste es eigentlich gemessen werden?
  \item Warum führt eine reine ``Reichweite''-Optimierung in B2B fast immer zu schlechten Ergebnissen?
\end{enumerate}

\ytquery{Performance Marketing Funnel Reichweite Traffic Lead Conversion Umsatz}

\antwortlinien{10}

\clearpage

% --- AUFGABE 2 -----------------------------------------------------------
\aufgabenkopf{2}{Zehn Performance-Marketing-KPIs}{\nivLi}{8\,min}

\noindent Aus dem Lehrskript Tag 9: \emph{Impressionen \textperiodcentered{} Klickrate \textperiodcentered{} Kosten pro Klick (CPC) \textperiodcentered{} Conversion Rate \textperiodcentered{} Kosten pro Lead (CPL) \textperiodcentered{} Leadqualität \textperiodcentered{} Abschlussquote \textperiodcentered{} Customer Acquisition Cost (CAC) \textperiodcentered{} Return on Ad Spend (ROAS) \textperiodcentered{} Deckungsbeitrag pro Kampagnenquelle}.

\medskip
\noindent\textbf{Aufgabe:}

\begin{enumerate}
  \item Ordnen Sie jede KPI einer Ebene zu: \emph{Anzeige \textperiodcentered{} Klick/Landing \textperiodcentered{} Lead \textperiodcentered{} Vertrieb \textperiodcentered{} Wirtschaftlichkeit}.
  \item Welche KPI ist der \emph{strategisch wichtigste Frühindikator} für die Wirtschaftlichkeit einer Kampagne?
  \item Welche KPI wird im Mittelstand am häufigsten \emph{nicht berichtet} \textendash{} und welche Folgen hat das?
\end{enumerate}

\ytquery{Performance Marketing KPI CPC CPL CAC ROAS Conversion Funnel}

\antwortlinien{14}

\clearpage

% --- AUFGABE 3 -----------------------------------------------------------
\aufgabenkopf{3}{Anzeigenkanäle nach Suchintention}{\nivLi}{6\,min}

\noindent Suchintentionen werden typischerweise in drei Stufen unterteilt: \emph{Problemrecherche \textperiodcentered{} Lösungsvergleich \textperiodcentered{} konkrete Kaufabsicht}. Anzeigenkanäle bedienen unterschiedliche Stufen.

\medskip
\noindent\textbf{Aufgabe:}

\begin{enumerate}
  \item Ordnen Sie diese Kanäle den drei Stufen zu (Mehrfachzuordnung möglich): \emph{Google Ads (Suchanfragen) \textperiodcentered{} LinkedIn Ads (Targeting nach Rolle) \textperiodcentered{} Meta Ads (Interessen) \textperiodcentered{} Display-Anzeigen \textperiodcentered{} Retargeting \textperiodcentered{} Plattformanzeigen \textperiodcentered{} Partnerkampagnen}.
  \item Welcher Kanal ist für die Phase \emph{konkrete Kaufabsicht} typischerweise am wirksamsten?
  \item Welche zwei Kanäle sind im B2B-Mittelstand für \emph{Buying-Center-Ansprache} besonders relevant?
\end{enumerate}

\ytquery{Performance Marketing Google Ads LinkedIn Suchintention Kaufabsicht B2B}

\antwortlinien{12}

\clearpage

\section*{Niveau L2 \textendash{} Anwendung \& Transfer}
\addcontentsline{toc}{section}{L2 -- Anwendung}

% --- AUFGABE 4 -----------------------------------------------------------
\aufgabenkopf{4}{SolarOffice GmbH \textendash{} Performance-Marketing-Funnel}{\nivLii}{25\,min}

\begin{infobox}[Fallbeschreibung]
``SolarOffice GmbH'' verkauft Photovoltaiklösungen für Gewerbeimmobilien. Anfragen entstehen bisher über Empfehlungen und regionale Messen. Die Geschäftsführung möchte über digitale Anzeigen \emph{mehr qualifizierte Beratungstermine} gewinnen.

\smallskip
\textbf{Gegeben:}
\begin{itemize}
  \item Zielgruppe: Eigentümer und Geschäftsführer kleiner und mittlerer Unternehmen.
  \item Produkt: erklärungsbedürftig, Investitionssumme oft $>$\,50.000\,\euro.
  \item Kaufmotiv: Energiekosten senken, Nachhaltigkeit, Fördermöglichkeiten.
  \item Verkaufszyklus: Wochen bis Monate.
  \item Vertriebsteam begrenzt, will keine unqualifizierten Anfragen.
  \item Budget erste Kampagne: 25.000\,\euro.
\end{itemize}
\end{infobox}

\noindent\textbf{Arbeitsauftrag:}

\begin{enumerate}
  \item Beschreiben Sie das \textbf{übergeordnete Vertriebsziel} der Kampagne in 2 Sätzen (nicht ``mehr Leads'').
  \item Benennen Sie mindestens \textbf{drei mögliche Zielgruppen} (z.\,B.\ produzierendes Gewerbe, Immobilieneigentümer, Logistikbetriebe).
  \item Formulieren Sie \textbf{drei Kundenprobleme}, die in Anzeigen angesprochen werden könnten.
  \item Skizzieren Sie einen einfachen \textbf{Funnel} von der Anzeige bis zum Beratungstermin (5\,Stufen).
  \item Legen Sie \textbf{drei Kennzahlen} fest, mit denen der Kampagnenerfolg gemessen wird \textendash{} nicht ``Klicks''.
\end{enumerate}

\ytquery{Performance Marketing B2B PV Solar Gewerbe Zielgruppe Funnel Kampagne}

\antwortlinien{20}

\clearpage

% --- AUFGABE 5 -----------------------------------------------------------
\aufgabenkopf{5}{SolarOffice \textendash{} Anzeigenkanäle \& Landing Page}{\nivLii}{30\,min}

\begin{infobox}[Anzeigenkanäle + Landing-Page-Zustand]
\textbf{Anzeigenkanäle:}
\begin{itemize}
  \item \textbf{A:} Google Ads für ``Photovoltaik Gewerbe Kosten''.
  \item \textbf{B:} LinkedIn Ads für Geschäftsführer und Immobilienverantwortliche.
  \item \textbf{C:} Display-Anzeigen auf allgemeinen Nachrichtenseiten.
  \item \textbf{D:} Retargeting für Websitebesucher.
  \item \textbf{E:} Plattformanzeigen auf Energie-/Nachhaltigkeitsportalen.
\end{itemize}

\smallskip
\textbf{Aktuelle Landing Page:}
\begin{itemize}
  \item Überschrift: ``Willkommen bei SolarOffice''.
  \item Viele allgemeine Informationen zum Unternehmen.
  \item Kein klarer Einstieg für Gewerbekunden.
  \item Formular mit 12 Pflichtfeldern.
  \item Keine Beispielrechnung. Keine Referenzen.
  \item CTA: ``Absenden''.
\end{itemize}

\smallskip
\textbf{Rahmen:} 80 qualifizierte Beratungstermine in 6\,Monaten \textperiodcentered{} max.\ 20 neue Leads pro Woche bearbeitbar \textperiodcentered{} erste Ergebnisse nach 8\,Wochen \textperiodcentered{} Marke seriös und beratungsorientiert.
\end{infobox}

\noindent\textbf{Arbeitsauftrag:}

\begin{enumerate}
  \item Bewerten Sie die 5 Anzeigenkanäle nach \emph{Zielgruppenpassung, Kostenrisiko, Leadqualität, Geschwindigkeit}.
  \item Wählen Sie max.\ \textbf{drei Kanäle} für die erste Kampagnenphase.
  \item Benennen Sie \textbf{fünf Schwächen} der aktuellen Landing Page.
  \item Entwickeln Sie \textbf{fünf konkrete Verbesserungen} für die Landing Page.
  \item Entscheiden Sie: Mehr Budget zuerst in Anzeigen oder in Landing Page? Begründen Sie in 3\,Sätzen.
\end{enumerate}

\ytquery{Performance Marketing Landing Page B2B Solar Gewerbe Google LinkedIn}

\antwortlinien{22}

\clearpage

% --- AUFGABE 6 -----------------------------------------------------------
\aufgabenkopf{6}{Elemente einer wirksamen Landing Page}{\nivLii}{15\,min}

\noindent Aus dem Lehrskript: Elemente einer wirksamen Landing Page: \emph{klare Überschrift \textperiodcentered{} Nutzenversprechen \textperiodcentered{} Zielgruppenansprache \textperiodcentered{} Vertrauenselemente / Referenzen \textperiodcentered{} Formular \textperiodcentered{} CTA \textperiodcentered{} visuelle Führung \textperiodcentered{} mobile Nutzbarkeit}.

\medskip
\noindent\textbf{Arbeitsauftrag:}

\begin{enumerate}
  \item Beschreiben Sie für \emph{jedes Element} in einem Halbsatz: Welche Rolle spielt es für die Conversion?
  \item Welche zwei Elemente werden im Mittelstand am häufigsten \emph{schwach} umgesetzt?
  \item Formulieren Sie für SolarOffice eine \textbf{starke Überschrift} und ein \textbf{Nutzenversprechen} (jeweils max.\ 2 Sätze).
  \item Welche \textbf{drei Vertrauenselemente} sind im B2B-Gewerbekundengeschäft besonders wirksam?
\end{enumerate}

\ytquery{Landing Page B2B Überschrift Nutzenversprechen Vertrauen Conversion}

\antwortlinien{16}

\clearpage

% --- AUFGABE 7 -----------------------------------------------------------
\aufgabenkopf{7}{DataSecure Solutions \textendash{} Performance-Marketing-Kampagne}{\nivLii}{30\,min}

\begin{infobox}[Fallstudie: DataSecure Solutions]
``DataSecure Solutions'' ist mittelständischer Anbieter von Cybersecurity-Lösungen für Unternehmen mit 100\,--\,1.000 Mitarbeitenden. Vertrieb bisher über Empfehlungen, Messen, Netzwerke. Markt wettbewerbsintensiver $\rightarrow$ digitale Leads gewünscht.

\smallskip
\textbf{Herausforderung:} Nicht möglichst viele Leads, sondern \emph{Entscheider mit echtem Bedarf}. Cybersecurity-Lösungen erklärungsbedürftig, vertrauenssensibel, oft mit mehreren Entscheidern.

\smallskip
\textbf{Rahmen:} 16\,Mio.\,\euro{} Umsatz \textperiodcentered{} Vertrieb 8 / Marketing 2 \textperiodcentered{} 150.000\,\euro{} für 6\,Monate \textperiodcentered{} Ziel: 180 qualifizierte Leads, 60 Demo-Termine, 20 Angebote \textperiodcentered{} Zielgruppen: IT-Leiter, GF, Datenschutz, Compliance \textperiodcentered{} Verkaufszyklus 3\,--\,6\,Monate \textperiodcentered{} Website bisher conversion-schwach.

\smallskip
\textbf{Kampagnenansätze:}
\begin{itemize}
  \item \textbf{A:} Google Ads für akute Suchanfragen ``Cybersecurity Beratung Mittelstand''.
  \item \textbf{B:} LinkedIn Ads mit Whitepaper für IT-Leiter und GF.
  \item \textbf{C:} Webinar-Kampagne ``Cyberrisiken im Mittelstand reduzieren''.
  \item \textbf{D:} Retargeting für Websitebesucher.
  \item \textbf{E:} Display-Kampagne mit breiter Reichweite.
\end{itemize}
\end{infobox}

\noindent\textbf{Arbeitsauftrag:}

\begin{enumerate}
  \item Analysieren Sie die \textbf{Zielgruppen + Informationsbedürfnisse} \textendash{} 4 Punkte.
  \item Entwickeln Sie eine \textbf{strategische Anzeigenplanung für 6 Monate}.
  \item Priorisieren Sie die Kampagnenansätze nach: Zielgruppenpassung, Leadqualität, Kosten, Geschwindigkeit, Vertriebswirkung.
  \item Entwickeln Sie eine \textbf{Landing Page Struktur} für eine Demo-Anfrage.
  \item Definieren Sie \textbf{Conversion-Ziele + KPIs} entlang des Funnels.
\end{enumerate}

\ytquery{Cybersecurity B2B Performance Marketing LinkedIn Webinar Landing Page Demo}

\antwortlinien{24}

\clearpage

% --- AUFGABE 8 -----------------------------------------------------------
\aufgabenkopf{8}{Funnel-Verluste analysieren}{\nivLii}{20\,min}

\noindent Beispielzahlen einer 30-Tage-Kampagne (fiktiv für eine B2B-Kampagne):

\medskip
\noindent\textbf{Funnel-Daten:}

\smallskip
\begin{tabularx}{\linewidth}{@{}lc@{}}
\toprule
\textbf{Stufe} & \textbf{Anzahl} \\
\midrule
Impressionen        & 800.000 \\
Klicks              & 24.000 (3\,\% Klickrate) \\
Landing-Page-Besuche & 22.000 (Tracking-Abweichung) \\
Formular-Aufrufe    & 4.400 (20\,\%) \\
Formular-Absendungen & 880 (20\,\%) \\
Marketing-Qualified Leads & 660 (75\,\%) \\
Sales-Accepted Leads & 132 (20\,\%) \\
Sales-Qualified Leads & 79 (60\,\%) \\
Angebote            & 32 (40\,\%) \\
Abschlüsse          & 6 (19\,\%) \\
\bottomrule
\end{tabularx}

\medskip
\noindent\textbf{Aufgabe:}

\begin{enumerate}
  \item Identifizieren Sie die \textbf{drei größten Conversion-Verluste} im Funnel \textendash{} \emph{absolut} und \emph{prozentual} betrachtet.
  \item Welcher Verlust ist primär ein \emph{Anzeigen}-Problem, welcher ein \emph{Landing-Page}-Problem, welcher ein \emph{Vertriebs}-Problem?
  \item Berechnen Sie \emph{Kosten pro SQL} bei einem Budget von 75.000\,\euro.
  \item Entwickeln Sie für jeden der drei Verluste eine \emph{konkrete Optimierungsmaßnahme}.
\end{enumerate}

\ytquery{Performance Marketing Funnel Conversion Verluste Optimierung B2B}

\antwortlinien{18}

\clearpage

\section*{Niveau L3 \textendash{} Management \& Entscheidungsarchitektur}
\addcontentsline{toc}{section}{L3 -- Management}

% --- AUFGABE 9 -----------------------------------------------------------
\aufgabenkopf{9}{DataSecure Solutions \textendash{} 6-Monats-Kampagnenentscheidung}{\nivLiii}{45\,min}

\noindent Sie sind \emph{Chief Revenue Officer} bei DataSecure Solutions. 150.000\,\euro{} Budget, 6\,Monate, Ziel: 180 qualifizierte Leads, 60 Demos, 20 Angebote.

\medskip
\noindent\textbf{Arbeitsauftrag:}

\begin{enumerate}
  \item Treffen Sie eine \textbf{priorisierte Kanalentscheidung} (max.\ 4 Kanäle aus A--E + Retargeting).
  \item Verteilen Sie das Budget von 150.000\,\euro{} auf Kanäle, Landing Page, Content, Tracking, Tests.
  \item Entwickeln Sie eine \textbf{Lead-Qualifizierungslogik} (MQL/SAL/SQL Definitionen + Übergaberegeln + SLA).
  \item Definieren Sie \textbf{sechs KPIs} entlang des Funnels (von Klick bis Umsatz).
  \item Beschreiben Sie die \textbf{Roadmap} in drei Phasen (jeweils 2\,Monate).
  \item Treffen Sie eine \textbf{Entscheidung gegen} einen reichweitenstarken Kanal \textendash{} aus Lead-Qualitäts- oder Vertriebskapazitätsgründen.
  \item Treffen Sie mindestens eine Entscheidung unter Unsicherheit mit klarer Annahme und Lernindikator.
\end{enumerate}

\ytquery{DataSecure Performance Marketing CRO Budget Kanal Lead Qualifizierung Senior}

\antwortlinien{36}

\clearpage

% --- AUFGABE 10 ----------------------------------------------------------
\aufgabenkopf{10}{Leadmenge vs.\ Leadqualität}{\nivLiii}{30\,min}

\noindent\textbf{Diskussionsaufgabe.}

\noindent Performance-Marketing-Reports zeigen typischerweise \emph{Leadmenge} (``500 Leads generiert!''). Der Vertrieb sieht aber, dass nur 40 davon ernsthafte Interessenten sind. Wo soll Marketing seinen Erfolg messen?

\medskip
\noindent\textbf{Arbeitsauftrag:}

\begin{enumerate}
  \item Beschreiben Sie in 3\,--\,4 Sätzen den Konflikt zwischen ``Leadmenge'' und ``Leadqualität'' aus Sicht eines Performance-Marketing-Managers.
  \item Entwickeln Sie eine \textbf{gemeinsame Erfolgs-KPI} für Marketing und Vertrieb, die strukturell beide Seiten ausrichtet.
  \item Welche Anreizgestaltung würde Performance-Marketing strukturell auf \emph{Qualität} statt Menge ausrichten?
  \item Welche Kennzahl entlarvt sofort eine ``Lead-Müll-Kampagne''?
  \item Welche Kanäle sind besonders anfällig für ``Lead-Müll'' \textendash{} und welche besonders gut für Qualität?
\end{enumerate}

\ytquery{Lead Menge Qualität Performance Marketing CAC ROAS Konflikt B2B}

\antwortlinien{20}

\clearpage

% --- AUFGABE 11 ----------------------------------------------------------
\aufgabenkopf{11}{Transferprojekt \textendash{} Performance-Marketing-Architektur}{\nivLiii}{45\,min}

\begin{infobox}[Rolle und Ausgangslage]
\textbf{Ihre Rolle:} Chief Revenue Officer / Head of Performance Marketing.

\smallskip
\textbf{Ausgangslage:} Mittelständisches B2B-Unternehmen will digitale Leadgenerierung professionalisieren. Anzeigen wurden bereits geschaltet, jedoch ohne klare Zielgruppenlogik, ohne systematische Landing-Page-Optimierung und ohne einheitliche Leadqualitäts-Bewertung. Marketing berichtet Klicks und Leads, Vertrieb bewertet Gespräche und Abschlüsse, Controlling betrachtet Kosten. Eine gemeinsame Performance-Architektur fehlt.

\smallskip
\textbf{Erwartung GF:} steuerbares System, das qualifizierte Nachfrage erzeugt, Vertriebskapazitäten sinnvoll nutzt und Marketingbudgets wirtschaftlich rechtfertigt.

\smallskip
\textbf{Rahmen:} Budget begrenzt \textperiodcentered{} Wettbewerbsdruck \textperiodcentered{} Vertriebskapazität begrenzt \textperiodcentered{} Zielgruppendaten unvollständig \textperiodcentered{} Marketing/Vertrieb mit anderen KPIs \textperiodcentered{} Landing Pages bisher schwach \textperiodcentered{} GF erwartet Entscheidungen trotz unsicherer Datenlage.
\end{infobox}

\noindent\textbf{Arbeitsauftrag:} Erarbeiten Sie schriftlich:

\begin{enumerate}
  \item \textbf{Zielbild} der Performance-Marketing-Architektur (5\,Sätze).
  \item Eine \textbf{Zielgruppen- und Buying-Center-Analyse}.
  \item Eine \textbf{Kampagnenstrategie} mit priorisierten Anzeigenkanälen.
  \item Eine \textbf{Botschaftenlogik} je Zielgruppe und Funnel-Phase.
  \item Eine \textbf{Landing-Page-Architektur} für qualifizierte Conversion.
  \item Eine \textbf{Leadqualifizierungs-Logik} (MQL/SAL/SQL-Definitionen, Übergaberegeln, SLA).
  \item Eine \textbf{KPI-Architektur} von Anzeige bis Umsatz.
  \item Eine \textbf{Budget-Logik} (Testen, Skalieren, Stoppen, Optimieren).
  \item Eine \textbf{Governance-Struktur} (Marketing, Vertrieb, Controlling, GF).
  \item \textbf{Drei Managemententscheidungen} unter Unsicherheit.
\end{enumerate}

\noindent\textbf{Zusatzanforderung:} Treffen Sie mindestens eine bewusste Entscheidung gegen einen Kanal / eine Kampagne, die viele Leads verspricht, aber voraussichtlich geringe Leadqualität, hohe Vertriebslast oder schlechte Wirtschaftlichkeit erzeugt.

\ytquery{Performance Marketing Architektur B2B CRO Conversion Lead Qualität Governance}

\antwortlinien{40}

\clearpage

\section*{Abschluss}
\thispagestyle{plain}

\noindent Sie haben alle elf Aufgaben \textendash{} und damit Tag 9 des Kurses 22 \textendash{} abgeschlossen.

\medskip
\begin{infobox}[Was Sie jetzt können sollten]
\begin{itemize}
  \item Performance-Marketing als Vertriebsprozess statt als Werbemaßnahme verstehen.
  \item Reichweite / Traffic / Leads / Conversion / Umsatz als Funnel-Logik trennen.
  \item Zehn KPIs entlang des Funnels einordnen und priorisieren.
  \item Anzeigenkanäle nach Suchintention und Zielgruppe einsetzen.
  \item Eine wirksame Landing Page strukturieren und Schwächen erkennen.
  \item Den strukturellen Konflikt Leadmenge vs.\ Leadqualität mit gemeinsamen KPIs adressieren.
  \item Eine Performance-Marketing-Architektur als Investitionsentscheidung denken.
\end{itemize}
\end{infobox}

\vspace{1cm}
\noindent{\color{lpAccent}\rule{\linewidth}{0.6pt}}\\[4pt]
{\footnotesize\sffamily\color{lpGray}Lernplatt \textperiodcentered{} by Can Siebert \textperiodcentered{} Kurs 22 (Bo 143) \textperiodcentered{} Tag 9 \textperiodcentered{} Aufgabenheft \textperiodcentered{} \textcopyright{} 2026}

\end{document}
